
%\needspace{4\baselineskip}
\color{uniblue}{\subsection{Логика отключения выключателя стороны НН трансформатора (ЛО НН)}}
\color{black}

\begin{enumerate}

%%%%%%%%%%%%%%%%%%%%%%%%%%%%%%%%%%%%%%%%%%%%%%%%%%%%%%%%%%%%%%%%%%%%% ОБЩИЕ СВЕДЕНИЯ %%%%%%%%%%%%%%%%%%%%%%%%%%%%%%%%%%%%%%%%%%%%%%%%%%%%%%%%%%%%%%

\item Общие сведения

\begin{enumerate}

\item
Логика отключения выключателя стороны НН трансформатора служит для формирования команды отключения в схему выключателя ввода стороны НН трансформатора, команды запрета АПВ в схему АПВ выключателя ввода секции НН и команды запрета АВР в схему АВР секций НН.

\item
В устройстве логика выключателя стороны НН трансформатора представлена функциональным блоком «ЛО НН». Функциональный блок «ЛО НН» состоит непосредственно из логики отключения («ЛО»), логики запрета АПВ («ЗАПВ») и логики запрета АВР («ЗАВР»).

% ЗДЕСЬ ВЕТВЛЕНИЕ, если исполнение Т2 то вставляем следующий абзац:
\ifboolexpr{bool {isT2}}
{
    \item
    В составе устройства предусматриваются две функциональных блока «ЛО НН1» и «ЛО НН2» (для стороны НН1 и НН2 соответственно), ниже приводится общее описание работы функционального блока «ЛО НН». 
}
{}

\item
ЛО НН может быть выведена из работы оперативно путем активации сигнала «ОВ~ЛО~НН».
\item
В таблице \ref{lonn:tbl_lov} приведены параметры, необходимые для настройки ЛО НН.

\end{enumerate}

%%%%%%%%%%%%%%%%%%%%%%%%%%%%%%%%%%%%%%%%%%%%%%%%%%%%%%%%%%%%%%%% ЛОГИКА ОТКЛЮЧЕНИЯ ВЫКЛЮЧАТЕЛЯ НН %%%%%%%%%%%%%%%%%%%%%%%%%%%%%%%%%%%%%%%%%%%%%%%%%%

\item Логика отключения выключателя ввода стороны НН (ЛО)

\begin{enumerate}

\item
Логика отключения в составе функционального блока «ЛО НН» представлена функцией «ЛО».

\item
Ввод функции в работу на этапе параметрирования устройства осуществляется программным переключателем SGF1 «Ввод\_функции». Информация о введенной функции отображается выходным сигналом «ЛО НН / ЛО: Ввод».

\item
Функция  может быть выведена из работы оперативно путем активации сигнала «ОВ ЛО НН / ЛО», что характеризуется активным состоянием сигнала «ЛО НН / ЛО: Оперативный вывод» на выходе функции.

\item
Функция принимает сигнал срабатывания от общей логики отключения трансформатора и формирует сигналы «ЛО НН / ЛО: Отключение» и «ЛО НН / ЛО: Отключить аварийно», сигнал «ЛО НН / ЛО: Отключить аварийно» является импульсным, время импульса определяется выдержкой времени таймера T1 (уставка «Тимп»).

\item
Алгоритм работы функции выполнен в соответствии с рисунком \ref{lonn:img_lov}.

\vspace{2mm}
\begin{figure}[H]
\centering
\includegraphics[width=0.9\textwidth,height=0.9\textheight,keepaspectratio]{\fbpath/36. ЛО В НН Т 35 Т/_latex/img/lov.pdf}
\captionof{figure}{Функциональная схема логики отключения выключателя}\label{lonn:img_lov}
\end{figure}

\end{enumerate}
%%%%%%%%%%%%%%%%%%%%%%%%%%%%%%%%%%%%%%%%%%%%%%%%%%%%%%%%%%%%%%%% ЛОГИКА ЗАПРЕТА АПВ НН %%%%%%%%%%%%%%%%%%%%%%%%%%%%%%%%%%%%%%%%%%%%%%%%%%

\item Логика запрета АПВ выключателя ввода стороны НН (ЗАПВ)

\begin{enumerate}

\item
Логика запрета АПВ в составе функционального блока «ЛО НН» представлена функцией «ЗАПВ».

\item
Ввод функции в работу на этапе параметрирования устройства осуществляется программным переключателем SGF1 «Ввод\_функции». Информация о введенной функции отображается выходным сигналом «ЛО НН / ЗАПВ: Ввод».

\item
Функция может быть выведена из работы оперативно путем активации сигнала «ОВ ЛО НН / ЗАПВ», что характеризуется активным состоянием сигнала «ЛО НН / ЗАПВ: Оперативный вывод» на выходе функции.

\item
Функция принимает сигнал срабатывания от общей логики отключения и формирует сигналы «ЛО НН / ЗАПВ: Запрет АПВ».

\item
Алгоритм работы функции выполнен в соответствии с рисунком \ref{lonn:img_zapv}.

\vspace{2mm}
\begin{figure}[H]
\centering
\includegraphics[width=0.9\textwidth,height=0.9\textheight,keepaspectratio]{\fbpath/36. ЛО В НН Т 35 Т/_latex/img/zapv.pdf}
\captionof{figure}{Функциональная схема логики запрета АПВ}\label{lonn:img_zapv}
\end{figure}

\end{enumerate}
%%%%%%%%%%%%%%%%%%%%%%%%%%%%%%%%%%%%%%%%%%%%%%%%%%%%%%%%%%%%%%%% ЛОГИКА ЗАПРЕТА АВР НН %%%%%%%%%%%%%%%%%%%%%%%%%%%%%%%%%%%%%%%%%%%%%%%%%%

\item Логика запрета АВР стороны НН (ЗАВР)

\begin{enumerate}

\item
Логика запрета АВР в составе функционального блока «ЛО НН» представлена функцией «ЗАВР».

\item
Ввод функции в работу осуществляется программным переключателем SGF1 «Ввод\_функции». Информация о введенной функции отображается выходным сигналом «ЛО НН / ЗАВР: Ввод».

\item
Функция может быть выведена из работы оперативно путем активации сигнала «ОВ ЛО НН / ЗАВР», что характеризуется активным состоянием сигнала «ЛО НН / ЗАВР: Оперативный вывод» на выходе функции.
\item
Функция принимает сигнал срабатывания от общей логики отключения и формирует сигналы «ЛО НН / ЗАВР: Запрет АВР».

\item
Алгоритм работы функции выполнен в соответствии с рисунком \ref{lonn:img_zavr}. 

\vspace{2mm}
\begin{figure}[H]
\centering

\includegraphics[width=0.9\textwidth,height=0.9\textheight,keepaspectratio]{\fbpath/36. ЛО В НН Т 35 Т/_latex/img/zavr.pdf}
\captionof{figure}{Функциональная схема логики запрета АВР}\label{lonn:img_zavr}
\end{figure}

\end{enumerate}

\end{enumerate}